% Define the document type and font size
\documentclass[12pt]{article}

\usepackage[spanish]{babel} % Package for the Spanish language
\usepackage[utf8]{inputenc} % Package for UTF-8 character encoding
\usepackage{geometry} % Package to configure document margins
\usepackage{graphicx} % Package to include graphics
\usepackage{listings} % Package to include source code in the document
\usepackage{xcolor} % Package to define colors
\usepackage{fancyhdr} % Package to customize headers and footers
\usepackage{amsmath} % Package for advanced mathematical symbols and environments
\usepackage{amssymb} % Package for additional mathematical symbols
\usepackage[hidelinks]{hyperref} % Package for hyperlinks in index and references (Keep this as the last package loaded)

\setlength{\parskip}{1em} % Space between paragraphs
\setlength{\parindent}{0.0in} % Indentation for paragraphs

% Definition of colors for source code
\definecolor{COLOR_COMMENTS}{HTML}{999999}
\definecolor{COLOR_KEYWORDS}{HTML}{37a5a6}
\definecolor{COLOR_IDENTIFIERS}{HTML}{ea76cb}
\definecolor{COLOR_STRINGS}{HTML}{ff7858}

% Configuration of the listings environment to display source code
\lstset{
frame=shadowbox,
aboveskip=3mm,
belowskip=3mm,
xleftmargin=10mm,
xrightmargin=10mm,
showstringspaces=false,
columns=flexible,
basicstyle={\footnotesize\ttfamily},
numbers=left,
numberstyle=\tiny\color{COLOR_COMMENTS},
keywordstyle=\color{COLOR_KEYWORDS},
stringstyle=\color{COLOR_STRINGS},
commentstyle=\color{COLOR_COMMENTS},
breaklines=true,
breakatwhitespace=true,
tabsize=4,
}

% Configuration of document margins
\geometry{
    a4paper,
    left=25mm,
    right=25mm,
    top=25mm,
    bottom=25mm
}
\graphicspath{{C:/Users/mishi/projects/LaTeX}} % Path for images

\title{Compresión de imágenes con SVD} % Document title
\author{María Laura Reina y Victoria Ruza} % Document author
\date{\today} % Document date

\begin{document}

    \begin{titlepage}
        \centering
        {\includegraphics[width=0.2\textwidth]{logoBlack}\par}
        {\large {Universidad Central de Venezuela}\par}
        {\large {Facultad de Ciencias}\par}
        {\large {Escuela de Computación}\par}
        {\large {Cálculo científico}\par}
        \vfill
        {\Huge {Compresión de imágenes con SVD}\par}
        \vspace{0.5cm}
        {\Large {Análisis del proyecto}\par}
        \vfill
        {\large María Laura Reina y Victoria Ruza\par}
        \vspace{0.5cm}
        {\large \today\par}
    \end{titlepage}

    \pagenumbering{gobble}
    \tableofcontents
    \newpage
    \pagenumbering{arabic}

    \section{Objetivos}
        \subsection{Objetivo General}
            \begin{itemize}
                \item Implementar y analizar la descomposición en valores singulares (SVD) como herramienta de compresión de imágenes, tomando en cuenta el balance entre calidad de la imagen y tamaño del archivo.
            \end{itemize}

        \subsection{Objetivos Específicos}
            \begin{itemize}
                \item Comprender los fundamentos matemáticos de la SVD y su aplicación en la compresión de imágenes.
                \item Desarrollar un entorno en Python capaz de procesar imagenes de alta resolución y transformarlas en matrices a las que es aplicable la SVD.
                \item Generar reconstrucciones dadas cantidades ascendentes de componentes principales $k /in \{5, 20, 50, 100, 200\}$.
                \item Analizar los resultados obtenidos haciendo una comparación de diferentes métricas, incluyendo la relación de almacenamiento, error relativo y energía capturada.
            \end{itemize}

    \newpage

    \section{Implementación Computacional}


    \section{Análisis Cuantitativo}
        \subsection{Tasa de Compresión}
            La tasa de compresión de datos se puede definir como una medida relativa que indica cuánto se ha reducido el tamaño de un archivo o conjunto de datos en comparación con su tamaño original. Se expresa comúnmente como un porcentaje o una relación entre el tamaño original y el tamaño comprimido. La fórmula general para calcular esta tasa es:
            \[
                \text{Tasa de compresión} = \frac{\text{Tamaño Original} - \text{Tamaño Comprimido}}{\text{Tamaño Original}} \times 100\%
            \]
            El tamaño original y comprimido serán medidos en datos con el propósito de simplificar el análisis abstrayendo el tamaño necesario para almacenar cada elemento de las matrices.

            La cantidad de datos necesarios para almacenar la imagen original es igual a $m \times n$, mientras que la cantidad de datos necesarios para almacenar la imagen comprimida es igual a $k \times (m + n + 1)$. Donde $m$ y $n$ son las dimensiones de la imagen original, y $k$ es el número de componentes principales o valores singulares utilizados en la reconstrucción. Por lo tanto, se puede obtener el tamaño original de la siguiente manera:
            \[
                \text{Tamaño Original} = m \times n
            \]



        \subsection{Análisis de Error}

        \subsection{Energía Capturada}

    \section*{Referencias}
        \begin{itemize}
            \item[[1]] 
        \end{itemize}

\end{document}